% Clase 13 --- Qué se está suponiendo al dibujar un circuito (§2.1 y §2.2).
% Las hipótesis de la "aproximación de circuito" son fáciles de usar sin notar
% que son hipótesis. El panel (a) las pone todas sobre el mismo dibujo que hizo
% el docente (batería, interruptor, resistencias) y el (b) muestra la
% consecuencia inmediata de que no se acumule carga: la ley de los nudos.
\begin{center}
\begin{tikzpicture}[scale=0.95]

  % =============== (a) el circuito y sus hipótesis ===============
  \begin{scope}
    \begin{scope}[line width=0.9pt, color=figblue]
      \draw (0,0) to[battery1, l=$\varepsilon$] (0,2.6);
      \draw (0,2.6) to[switch, l=interruptor] (3.2,2.6);
      \draw (3.2,2.6) to[R, l=$R_1$] (3.2,0);
      \draw (3.2,0) to[R, l_=$R_2$] (0,0);
    \end{scope}

    % corriente
    \draw[figvecres] (1.1,2.15) -- (2.1,2.15);
    \node[figlbl, figred, anchor=north] at (1.6,2.1) {$I$};

    % cable de sección variable: la misma corriente
    \draw[figgray, line width=0.5pt] (0.9,-0.16) -- (2.3,-0.16);
    \draw[figgray, line width=0.5pt] (0.9,0.16) -- (2.3,0.16);
    \node[figlblaux, anchor=north, align=center] at (1.6,-0.35)
      {aunque cambie la sección,\\[-0.2em] la corriente es la misma};

    \node[figlblaux, anchor=east, align=right] at (-0.9,1.3)
      {cable ideal:\\[-0.2em] resistencia\\[-0.2em] despreciable};

    \node[figlbl, anchor=north, align=center] at (1.6,-1.5)
      {(a) la aproximación de circuito};
  \end{scope}

  % =============== (b) el nudo ===============
  \begin{scope}[xshift=9.6cm, yshift=1.3cm]
    \begin{scope}[line width=0.9pt, color=figblue]
      \draw (-2.2,0) -- (0,0);
      \draw (0,0) -- (2.0,1.1);
      \draw (0,0) -- (2.0,-1.1);
      \fill (0,0) circle (2.2pt);
    \end{scope}
    \node[figlblaux, anchor=north] at (0,-0.2) {nudo};

    \draw[figvecres] (-1.9,0.28) -- (-0.9,0.28);
    \node[figlbl, figred, anchor=south] at (-1.4,0.32) {$I_1$};
    \draw[figvecres] (0.85,0.75) -- (1.75,1.24);
    \node[figlbl, figred, anchor=south east] at (1.5,1.2) {$I_2$};
    \draw[figvecres] (0.85,-0.75) -- (1.75,-1.24);
    \node[figlbl, figred, anchor=north east] at (1.5,-1.2) {$I_3$};

    \node[figlbl, figred, anchor=north] at (0,-1.95) {$I_1 = I_2 + I_3$};

    \node[figlbl, anchor=north, align=center] at (0,-2.8)
      {(b) ley de los nudos};
  \end{scope}

\end{tikzpicture}\\[0.5em]
{\small\itshape Las cuatro hipótesis se usan todo el tiempo sin nombrarlas.
Que el cable sea ideal significa que su resistencia es despreciable
\emph{frente a la de los resistores} ---no que sea nula: cuando no hay ningún
resistor, como en un cortocircuito, esa resistencia diminuta es la única que
queda y la corriente se vuelve enorme---. Que la corriente sea la misma a lo
largo de un tramo, incluso donde el cable se ensancha, no es geometría sino
conservación de la carga. Y que el circuito quede neutro en todas partes es lo
que hace que la corriente sea estacionaria: acumular carga en algún punto
costaría muchísima energía, así que el sistema se acomoda ---en un transitorio
de unos $10^{-9}$ s, imperceptible--- para que eso no pase. La ley de los nudos
del panel (b) es exactamente esa hipótesis aplicada a un punto de empalme: lo
que entra, sale.}
\end{center}
